\section{模型学习与基于模型的规划}

\subsection{从试错到想象}

\begin{figure}[H]
\centering
\includegraphics[width=0.95\textwidth]{slides-images/slide-030.jpg}
\caption{模型学习框架：学习物理世界的近似模型，在虚拟域中预测动作效果，指导真实世界中的动作决策\protect\footnotemark}
\end{figure}
\footnotetext{Slides 第31页。}

人类之所以能高效地与环境交互，是因为我们拥有\textbf{直觉物理模型}——能想象动作的后果。模型学习（Model Learning）试图赋予机器人类似的能力。

\begin{importantbox}{前向模型与逆规划}
\begin{itemize}
    \item \textbf{前向模型（Forward Model）}：给定当前状态 $s_t$ 和动作 $a_t$，预测下一状态 $s_{t+1}$
    \item \textbf{规划（Planning）}：前向模型的逆问题——给定当前状态 $s_t$ 和目标状态 $s^*$，求解动作序列 $\{a_t, a_{t+1}, \ldots\}$ 使预测轨迹尽可能接近目标
\end{itemize}
\end{importantbox}

规划的优化过程：
$$\min_{\{a_t\}} \sum_{t} \| f(s_t, a_t) - s^*_t \|^2$$
其中 $f$ 是学习到的前向模型。实际执行时采用\textbf{MPC（模型预测控制）}策略：只执行优化出的第一个动作，获取真实环境的新状态后重新规划，以应对模型误差。

\subsection{状态表征的选择}

状态表征对模型学习至关重要。不同任务需要不同粒度的表征：

\begin{figure}[H]
\centering
\includegraphics[width=0.95\textwidth]{slides-images/slide-034.jpg}
\caption{三种状态表征：像素动力学（左上）、关键点动力学（左下）、粒子动力学（右）\protect\footnotemark}
\end{figure}
\footnotetext{Slides 第35页。}

\begin{enumerate}
    \item \textbf{像素动力学}：直接以 2D 图像为状态，预测动作后图像如何变化。代表工作：Deep Visual Foresight
    \item \textbf{关键点动力学}：用 3D 关键点追踪物体，预测关键点的运动。适用于中等自由度的物体操控
    \item \textbf{粒子动力学}：用点集表示物体，预测粒子在动作下的运动。适用于高自由度物体（如颗粒物料、可形变物体）
\end{enumerate}

\subsection{实例：饺子机器人}

\begin{figure}[H]
\centering
\includegraphics[width=0.95\textwidth]{slides-images/slide-038.jpg}
\caption{饺子制作机器人：配备 15 种 3D 打印工具和 4 个 RGBD 相机，通过粒子动力学模型进行工具选择和动作规划\protect\footnotemark}
\end{figure}
\footnotetext{Slides 第39页。}

这项在 Stanford 完成的工作展示了模型学习的强大能力：机器人配备 15 种 3D 打印工具和 4 个 RGBD 相机，通过学习面团的粒子动力学模型，自主决策工具选择和动作规划，从一团面团制作出完整的饺子。

该系统的决策分为两个层次：
\begin{itemize}
    \item \textbf{高层决策}：选择使用哪种工具（分类器）
    \item \textbf{低层决策}：给定工具，规划具体的运动轨迹
\end{itemize}

\begin{knowledgebox}{学习模型 vs 物理仿真器}
实验表明，直接从真实世界交互数据学习的神经动力学模型，比精心调参的基于物理的仿真器（如 MPM 材料点方法）\textbf{预测精度更高}。这为``数据驱动的物理模型''提供了有力证据。
\end{knowledgebox}

\subsection{本章小结}

模型学习通过学习环境的前向预测模型，让机器人能``想象''动作的后果，从而进行高效规划。关键要素包括：合适的状态表征（像素/关键点/粒子）、准确的前向模型、以及 MPC 式的闭环执行。这种方法比纯强化学习更数据高效、更安全。

%% ========== Part 5: 模仿学习 ==========

